%% kapa: exam=1 round=22 year=2011 period=2 form=A subject=accounting subject_ko=회계학 seq=14 number=14
%% source: Q-Net 감정평가사 기출문제 (cst00309) artlSeq=5036790 file=extracted/1008/5036790/2011년 제22회 감정평가사 1차시험 기출문제/2교시.hwp
%% conversion: kordoc 4.24.51 -> tools/gen1.py (auto); review=proofread-gpt-astra; flags=none
\begin{question}{14}
다음 거래들 중에서 부채비율(=부채/자본)을 높이는 것만을 모두 고른 것은?
\begingroup
% 이 지문 상자는 중간에서 분할하지 않는다.
\renewtcolorbox{framedtext}{enhanced, unbreakable, colback=white, colframe=black,
  boxrule=0.4pt, arc=0pt, left=4pt, right=4pt, top=4pt, bottom=4pt}
\begin{framedtext}
ㄱ. 보유하고 있던 설비자산(장부금액 ￦10,000, 공정가치 ￦12,000)을 차량운반구(장부금액 ￦8,000, 공정가치 ￦10,000)와 교환하였는데, 상업적 실질이 결여되어 있다. \par
ㄴ. 보유중인 유휴자산(장부금액 ￦100,000)을 ￦40,000에 매각하고 현금 수취하였다. \par
ㄷ. 재고자산(장부금액 ￦30,000)을 재구매 조건으로 ￦50,000에 처분하였다. \par
ㄹ. 보관중인 무이자부 어음(액면금액 ￦10,000)을 금융기관에서 할인하여 현금 ￦7,000을 수취하였다. 본 거래는 금융자산의 제거조건을 충족한다. \par
ㅁ. 단기매매차익을 목적으로 시장성 있는 주식 10주를 주당 ￦12,000에 매입하고 매입수수료 ￦5,000과 함께 현금으로 지급하였다. \par
\end{framedtext}
\endgroup
\begin{choices}
  \item ㄷ
  \item ㄷ, ㄹ
  \item ㄱ, ㄴ, ㄷ, ㄹ
  \item ㄴ, ㄷ, ㄹ, ㅁ
  \item ㄱ, ㄴ, ㄷ, ㄹ, ㅁ
\end{choices}
\end{question}
\answer{4}
